\chapter{Binomial Identities}
\signature

\section{Two kinds of proof}

Recall $\binom{n}{k} = \frac{n!}{k!(n-k)!}$ counts the $k$-subsets of an
$n$-set. Binomial identities admit two complementary proof styles:
\emph{algebraic} (manipulate factorials, or compare coefficients in the
binomial theorem) and \emph{combinatorial} (``double counting'': count one
set in two ways). Whenever possible we give both --- the combinatorial proof
usually explains \emph{why} the identity holds.

\section{The fundamental identities}

\begin{theorem}[Symmetry]
$\displaystyle\binom{n}{k} = \binom{n}{n-k}$.
\end{theorem}

\begin{proof}
Algebraic: both equal $\frac{n!}{k!(n-k)!}$. Combinatorial: selecting $k$
elements is the same as rejecting $n - k$.
\end{proof}

\begin{theorem}[Pascal]
$\displaystyle\binom{n+1}{k} = \binom{n}{k-1} + \binom{n}{k}$.
\end{theorem}

\begin{proof}
Fix a distinguished element $x$ of an $(n+1)$-set: $k$-subsets containing
$x$ number $\binom{n}{k-1}$, those avoiding $x$ number $\binom{n}{k}$.
\end{proof}

Pascal's triangle hides many patterns: the hockey-stick sums below, the
Fibonacci numbers along shallow diagonals
($\sum_k \binom{n-k}{k} = F_{n+1}$), the parity fractal of Sierpi\'nski, and
the central column $\binom{2n}{n}$ governing Catalan numbers.

\begin{theorem}[Absorption]
$\displaystyle k\binom{n}{k} = n\binom{n-1}{k-1}$, equivalently
$\binom{n}{k} = \frac{n}{k}\binom{n-1}{k-1}$.
\end{theorem}

\begin{proof}
Count pairs (committee of $k$, its chair) from $n$ people in two ways:
choose the committee then the chair ($\binom{n}{k} \cdot k$), or the chair
then the rest ($n \cdot \binom{n-1}{k-1}$).
\end{proof}

\section{The binomial theorem and its harvest}

\begin{theorem}[Binomial theorem]
$\displaystyle (x+y)^{n} = \sum_{k=0}^{n} \binom{n}{k} x^{n-k} y^{k}$.
\end{theorem}

Substituting values of $x, y$ yields identities mechanically.

\begin{theorem}[Row sum]\label{thm:rowsum}
$\displaystyle\sum_{k=0}^{n} \binom{n}{k} = 2^{n}$.
\end{theorem}

\begin{proof}
Three proofs. (1) Set $x = y = 1$ in the binomial theorem.
(2) Combinatorial: both sides count subsets of an $n$-set, the left side by
size. (3) Induction via Pascal's identity.
\end{proof}

\begin{theorem}[Alternating sum]
$\displaystyle\sum_{k=0}^{n} (-1)^{k} \binom{n}{k} = 0$ for $n \geq 1$ ---
set $x = 1$, $y = -1$. Hence even-sized subsets equal odd-sized subsets in
number ($2^{n-1}$ each).
\end{theorem}

\begin{theorem}[Weighted sums]\label{thm:weighted}
\[
\sum_{k=0}^{n} 2^{k}\binom{n}{k} = 3^{n}, \qquad
\sum_{k=0}^{n} k\binom{n}{k} = n\,2^{n-1}, \qquad
\sum_{k=0}^{n} k^{2}\binom{n}{k} = n(n+1)\,2^{n-2}.
\]
\end{theorem}

\begin{proof}
First: $x=1, y=2$. Second: absorption gives
$\sum_k k\binom{n}{k} = n \sum_k \binom{n-1}{k-1} = n\,2^{n-1}$;
combinatorially, count (subset, distinguished member) pairs.
Third: write $k^2 = k(k-1) + k$ and apply absorption twice:
$\sum k(k-1)\binom{n}{k} = n(n-1)2^{n-2}$, then add $n\,2^{n-1}$ to get
$n(n-1)2^{n-2} + n\,2^{n-1} = n(n+1)2^{n-2}$.
\end{proof}

\section{Vandermonde and friends}

\begin{theorem}[Vandermonde]\label{thm:vdm}
$\displaystyle\binom{m+n}{r} = \sum_{k=0}^{r}\binom{m}{k}\binom{n}{r-k}$.
\end{theorem}

\begin{example}
A jury of $4$ from $6$ magistrates and $5$ citizens:
$\binom{11}{4} = 330$ directly; by cases on magistrates:
$\binom{6}{0}\binom{5}{4} + \binom{6}{1}\binom{5}{3} +
\binom{6}{2}\binom{5}{2} + \binom{6}{3}\binom{5}{1} +
\binom{6}{4}\binom{5}{0} = 5 + 60 + 150 + 100 + 15 = 330$. \checkmark
\end{example}

\begin{theorem}[Sum of squares]
$\displaystyle\sum_{k=0}^{n} \binom{n}{k}^{2} = \binom{2n}{n}$ ---
Vandermonde with $m = n = r$, using symmetry
$\binom{n}{n-k} = \binom{n}{k}$.
\end{theorem}

\begin{theorem}[Hockey stick]\label{thm:hockey2}
$\displaystyle\sum_{i=r}^{n} \binom{i}{r} = \binom{n+1}{r+1}$.
\end{theorem}

\begin{proof}
Classify $(r+1)$-subsets of $\{1, \ldots, n+1\}$ by largest element $i+1$:
the remaining $r$ elements come from $\{1,\dots,i\}$, in $\binom{i}{r}$
ways, for $i = r, \ldots, n$.
\end{proof}

\section{Extensions}

\begin{definition}[Generalised binomial coefficient]
For real $\alpha$ and $k \in \N$:
$\binom{\alpha}{k} = \frac{\alpha(\alpha-1)\cdots(\alpha-k+1)}{k!}$.
Newton's binomial series $(1+x)^{\alpha} = \sum_k \binom{\alpha}{k}x^{k}$
converges for $|x| < 1$; the \emph{negative-binomial} case gives
$\binom{-n}{k} = (-1)^{k}\binom{n+k-1}{k}$, linking to stars and bars.
\end{definition}

\begin{definition}[Catalan numbers]
$C_n = \dfrac{1}{n+1}\dbinom{2n}{n} = \dbinom{2n}{n} - \dbinom{2n}{n+1}$:
$1, 1, 2, 5, 14, 42, \ldots$ They count balanced parenthesisations,
monotonic lattice paths not crossing the diagonal, triangulations of convex
polygons, and binary trees.
\end{definition}

The multinomial coefficients $\binom{n}{n_1,\dots,n_m}$ extend the whole
theory from two blocks to $m$; and \emph{generating functions} --- viewing
$\sum_k \binom{n}{k}x^k = (1+x)^n$ as a polynomial identity --- turn
combinatorial convolutions into products of series, the method behind the
Vandermonde proof and far beyond.

\section{The combinatorial-proof template}

To prove $\mathrm{LHS} = \mathrm{RHS}$ combinatorially: (1) pose a concrete
counting question; (2) show the LHS answers it one way; (3) show the RHS
answers the same question another way; (4) conclude the counts are equal.
Practise the template on the worked example below.

\begin{example}[Portfolio choice]
Maeve must pick a team of $k$ analysts from $n$ candidates and name one team
lead, but she may also count by first naming the lead ($n$ ways) and then
completing the team ($\binom{n-1}{k-1}$ ways). The two counts,
$k\binom{n}{k}$ and $n\binom{n-1}{k-1}$, must agree: absorption, proved by
the template.
\end{example}

\section{Master cheat sheet}

\begin{center}
\renewcommand{\arraystretch}{1.6}
\begin{tabular}{ll}
\toprule
Symmetry & $\binom{n}{k} = \binom{n}{n-k}$\\
Pascal & $\binom{n+1}{k} = \binom{n}{k-1} + \binom{n}{k}$\\
Absorption & $k\binom{n}{k} = n\binom{n-1}{k-1}$\\
Row sum & $\sum_k \binom{n}{k} = 2^{n}$\\
Alternating & $\sum_k (-1)^{k}\binom{n}{k} = 0 \quad (n \geq 1)$\\
Weighted & $\sum_k k\binom{n}{k} = n2^{n-1}$, \quad
  $\sum_k k^{2}\binom{n}{k} = n(n+1)2^{n-2}$\\
Powers of 3 & $\sum_k 2^{k}\binom{n}{k} = 3^{n}$\\
Vandermonde & $\binom{m+n}{r} = \sum_k \binom{m}{k}\binom{n}{r-k}$\\
Squares & $\sum_k \binom{n}{k}^{2} = \binom{2n}{n}$\\
Hockey stick & $\sum_{i=r}^{n} \binom{i}{r} = \binom{n+1}{r+1}$\\
Catalan & $C_n = \frac{1}{n+1}\binom{2n}{n}$\\
\bottomrule
\end{tabular}
\end{center}

\begin{notebox}
Common pitfalls: forgetting that $\binom{n}{k} = 0$ for $k < 0$ or $k > n$;
shifting summation indices incorrectly; applying symmetry inside a sum
without adjusting the range; and assuming an identity for real upper index
when it only holds for integers.
\end{notebox}

\section{Summary}

Every identity in this chapter flows from three sources: Pascal's recurrence,
the binomial theorem with clever substitutions, and double counting. Master
the cheat sheet, but more importantly master the \emph{template}: a binomial
identity is an invitation to find the single question that both sides answer.

\section*{Exercises}
\addcontentsline{toc}{section}{Exercises}

\begin{exercise}{\easy}
Evaluate without a calculator: $\binom{10}{8}$, $\binom{12}{1}$,
$\binom{9}{4} + \binom{9}{5}$.
\end{exercise}

\begin{exercise}{\easy}
Use the row-sum and alternating-sum identities to find
$\sum_{k \text{ even}} \binom{n}{k}$ for $n \geq 1$.
\end{exercise}

\begin{exercise}{\easy}
Verify the hockey-stick identity
$\binom{2}{2} + \binom{3}{2} + \binom{4}{2} + \binom{5}{2} = \binom{6}{3}$
numerically.
\end{exercise}

\begin{exercise}{\medium}
Give a combinatorial proof of
$\binom{n}{2}\binom{n-2}{k-2} = \binom{n}{k}\binom{k}{2}$ \,($2 \le k \le
n$).
\end{exercise}

\begin{exercise}{\medium}
Evaluate $\sum_{k=0}^{n} 3^{k} \binom{n}{k}$ and
$\sum_{k=0}^{n} (-2)^{k} \binom{n}{k}$.
\end{exercise}

\begin{exercise}{\medium}
Prove $\sum_{k=1}^{n} k\binom{n}{k} = n2^{n-1}$ combinatorially, by counting
committees with a chair.
\end{exercise}

\begin{exercise}{\medium}
Using Vandermonde, evaluate $\sum_{k=0}^{m} \binom{m}{k}\binom{n}{p+k}$ as a
single binomial coefficient ($p + m \le n$).
\end{exercise}

\begin{exercise}{\medium}
Show that $\binom{2n}{n}$ is even for $n \geq 1$.
\emph{Hint: } $\binom{2n}{n} = \binom{2n-1}{n-1} + \binom{2n-1}{n}$ and
symmetry.
\end{exercise}

\begin{exercise}{\hard}
Prove the Fibonacci diagonal identity
$\sum_{k=0}^{\lfloor n/2 \rfloor} \binom{n-k}{k} = F_{n+1}$
(with $F_1 = F_2 = 1$), by induction using Pascal's identity.
\end{exercise}

\begin{exercise}{\hard}
Prove $\sum_{k=0}^{n} \binom{n}{k}^2 = \binom{2n}{n}$ combinatorially, by
counting $n$-subsets of a set of $n$ women and $n$ men.
\end{exercise}

\begin{exercise}{\hard}
Show $C_n = \binom{2n}{n} - \binom{2n}{n+1}$ counts lattice paths from
$(0,0)$ to $(n,n)$ (unit right/up steps) that never rise above the diagonal,
via the reflection principle.
\end{exercise}

\section*{Solutions to the Exercises}
\addcontentsline{toc}{section}{Solutions to the Exercises}

\begin{solution}{11.1}
$\binom{10}{8} = \binom{10}{2} = 45$; $\binom{12}{1} = 12$;
$\binom{9}{4} + \binom{9}{5} = \binom{10}{5} = 252$ (Pascal).
\end{solution}

\begin{solution}{11.2}
Adding $\sum_k \binom{n}{k} = 2^n$ and $\sum_k (-1)^k\binom{n}{k} = 0$
cancels odd $k$ and doubles even $k$:
$\sum_{k\text{ even}} \binom{n}{k} = 2^{n-1}$.
\end{solution}

\begin{solution}{11.3}
$1 + 3 + 6 + 10 = 20 = \binom{6}{3}$. \checkmark
\end{solution}

\begin{solution}{11.4}
Both sides count (committee of $k$ from $n$ people, with $2$ designated
co-chairs): left side picks the co-chairs first
($\binom{n}{2}$), then fills the committee ($\binom{n-2}{k-2}$); right side
picks the committee ($\binom{n}{k}$), then its co-chairs ($\binom{k}{2}$).
\end{solution}

\begin{solution}{11.5}
Binomial theorem with $(x,y) = (1,3)$: $4^{n}$. With $(1,-2)$:
$(-1)^{n}$.
\end{solution}

\begin{solution}{11.6}
Count pairs (committee $S$ of any size, chair $c \in S$) from $n$ people.
By size: choose size $k$, the committee, then the chair:
$\sum_k \binom{n}{k} k$. Chair first: $n$ choices, then any subset of the
remaining $n-1$ people joins: $n \, 2^{n-1}$. Equate.
\end{solution}

\begin{solution}{11.7}
By symmetry $\binom{m}{k} = \binom{m}{m-k}$, the sum is
$\sum_{k} \binom{m}{m-k}\binom{n}{p+k}$, a Vandermonde convolution picking
$m + p$ objects from $m + n$:
$\binom{m+n}{m+p}$.
\end{solution}

\begin{solution}{11.8}
Pascal: $\binom{2n}{n} = \binom{2n-1}{n-1} + \binom{2n-1}{n}$, and the two
terms are equal by symmetry ($(n-1) + n = 2n - 1$). Hence
$\binom{2n}{n} = 2\binom{2n-1}{n-1}$ is even.
\end{solution}

\begin{solution}{11.9}
Let $S_n = \sum_k \binom{n-k}{k}$. Check $S_1 = 1 = F_2$, $S_2 = 2 = F_3$.
For the step, Pascal gives
$\binom{n-k}{k} = \binom{n-1-k}{k} + \binom{n-1-k}{k-1}$;
summing over $k$, the first parts sum to $S_{n-1}$ and the second parts
re-index ($k \to k-1$) to $S_{n-2}$. So $S_n = S_{n-1} + S_{n-2}$ ---
the Fibonacci recurrence with Fibonacci initial values, hence
$S_n = F_{n+1}$ for all $n$.
\end{solution}

\begin{solution}{11.10}
Choosing $n$ people from $n$ women + $n$ men splits by the number $k$ of
women chosen: $\binom{n}{k}$ ways for the women and
$\binom{n}{n-k} = \binom{n}{k}$ for the men, giving
$\sum_k \binom{n}{k}^2$. Direct count: $\binom{2n}{n}$. Equate.
\end{solution}

\begin{solution}{11.11}
All paths to $(n,n)$: $\binom{2n}{n}$. A ``bad'' path touches the line
$y = x + 1$; reflecting its prefix up to the first touch across that line
bijects bad paths with paths from $(-1, 1)$ to $(n,n)$, which number
$\binom{2n}{n+1}$. Good paths: $\binom{2n}{n} - \binom{2n}{n+1} = C_n$.
\end{solution}
